In this section we introduce an extension to the semantics of CEL, namely we introduce the \textit{joins} operator to combine complex events that agree on their shared variables. 
We then extend the formal computational model for evaluating CEL formulas and prove its correctness. 

Given two complex events $C_1 = (i_1, j_1, \mu_1)$ and $C_2 = (i_2, j_2, \mu_2)$, we say that $C_1$ and $C_2$ are \emph{compatible}, denoted by $C_1 \sim C_2$, if $\mu_1(X) = \mu_2(X)$ for every $X \in \text{dom}(C_1) \cap \text{dom}(C_2)$. We present then the \emph{joins operator} for CEL with the following definition, where $\text{dom}(\varphi_1 \joins \varphi_2) = \text{dom}(\varphi_1) \cup \text{dom}(\varphi_2)$:
\[ 
\begin{aligned}
	\sem{\varphi_1 \joins \varphi_2}(\Stream)
	&= \{C_1 \cup C_2 \mid C_1 \in \sem{\varphi_1}(\Stream) 
	\land C_2 \in \sem{\varphi_2}(\Stream) \\
	&\quad \land \forall X \in \text{dom}(C_1) \cap \text{dom}(C_2).\ C_1(X) = C_2(X)\}
\end{aligned}
\]

\begin{example} Suppose we have sensors in a park monitoring humidity ($H$), temperature ($T$), and smoke ($S$). Consider a stream $\Stream = e_1 e_2 e_3$ where the event types are $\type(e_1) = H$, $\type(e_2) = T$, and $\type(e_3) = S$. 
	
	Let $\varphi_1 = H ; T$ be a formula detecting a humidity reading followed by a temperature reading, with $\text{dom}(\varphi_1) = \{H, T\}$, and $\varphi_2 = T ; S$ be a formula detecting a temperature reading followed by a smoke reading, with $\text{dom}(\varphi_2) = \{T, S\}$.
	\begin{itemize}
		\item The set $\sem{\varphi_1}(\Stream)$ contains a complex event $C_1$ with $\cinterval(C_1) = [1..2]$, $C_1(H) = \{1\}$, and $C_1(T) = \{2\}$.
		\item The set $\sem{\varphi_2}(\Stream)$ contains a complex event $C_2$ with $\cinterval(C_2) = [2..3]$, $C_2(T) = \{2\}$, and $C_2(S) = \{3\}$.
	\end{itemize}
	
	The shared domain between both events is $\text{dom}(C_1) \cap \text{dom}(C_2) = \{T\}$. Since $C_1(T) = C_2(T) = \{2\}$, the compatibility condition holds, and the complex event $C = C_1 \cup C_2$ is in the result of $\sem{\varphi_1 \joins \varphi_2}(\Stream)$ with $\cinterval(C) = [1..3]$, $C(H) = \{1\}$, $C(T) = \{2\}$, and $C(S) = \{3\}$. This resulting complex event uses the shared temperature reading at position $2$ as a bridge connecting the initial humidity event to the subsequent smoke event.
\end{example}

We denote by $\celplus{\joinssimb}$ the \emph{extension of CEL} with the joins operation. We can prove the following result:

\begin{theorem} \label{theo:joins}
	For every $\celplus{\joinssimb}$ formula $\varphi$ there exists a CEA $\cA_\varphi$ such that $\sem{\varphi}(\Stream) = \sem{\cA_\varphi}(\Stream)$ for every stream $\Stream$.
\end{theorem}

\begin{proof}
	We already know that for every (standard) CEL operator we can construct an equivalent CEA. So, we concentrate next in the construction of a CEA for the joins operator. 
	Let $\varphi_1$ and $\varphi_2$ be formulas in CEL extended with the joins operator. We show by induction that there exists a CEA $\cAjoins$ such that:
	\[
	\sem{\varphi_1 \joins \varphi_2}(\Stream) = \sem{\cAjoins}(\Stream)
	\]
	Assume then that there exist CEAs that satisfy the previous property for $\varphi_1$ and $\varphi_2$, called $\cA_{\varphi_1} = (Q_1, \SETP_1, \SETX_1, \Delta_1, q_0, F_1)$ and $\cA_{\varphi_2} = (Q_2, \SETP_2, \SETX_2, \Delta_2, p_0, F_2)$, respectively. 
	Then the construction for the joins operator is as follows. Let $\bot$ and $\top$ be two fresh state symbols representing that an automaton has not yet started or has already finished, respectively.
	Let $\cAjoins$ be a CEA where $\cAjoins = (\Qjoins, \Pjoins, \Xjoins, \Djoins, (\bot, \bot), \Fjoins)$:
	\[
	\begin{aligned}[t]
		& \Qjoins = (Q_1 \uplus \{\bot, \top\}) \times (Q_2 \uplus \{\bot, \top\}) \\
		& \Pjoins = \SETP_1 \cup \SETP_2 \cup \{ P_1 \land P_2 \mid P_1 \in \SETP_1 \land P_2 \in \SETP_2 \} \cup \{\text{true}\} \hspace*{120px}\\
		& \Xjoins  = \SETX_1 \cup \SETX_2 \\
		& (\bot, \bot) \text{ the initial state of }\cAjoins.\\
		& \Fjoins = \{(\top, \top)\} \text{ the final state of } \cAjoins.
	\end{aligned}
	\]
	The transition relation $\Djoins$ is formally defined as:
	\[
	\begin{aligned}[t]
		\Djoins & = \{((\bot, r), \epsilon, (q_0, r)) \mid r \in Q_2 \uplus \{\bot, \top\}\} \\
		& \uplus \{((r, \bot), \epsilon, (r, p_0)) \mid r \in Q_1 \uplus \{\bot, \top\}\} \\
		& \uplus \{((q, r), \epsilon, (\top, r)) \mid q \in F_1 \land r \in Q_2 \uplus \{\bot, \top\}\} \\
		& \uplus \{((r, p), \epsilon, (r, \top)) \mid r \in Q_1 \uplus \{\bot, \top\} \land p \in F_2\} \\
		& \uplus \{((q, r), \epsilon, (q', r)) \mid (q, \epsilon, q') \in \Delta_1 \land r \in Q_2 \uplus \{\bot, \top\}\} \\
		& \uplus \{((r, p), \epsilon, (r, p')) \mid r \in Q_1 \uplus \{\bot, \top\} \land (p, \epsilon, p') \in \Delta_2\} \\
		& \uplus \{((q, r), P_1, L_1, (q', r)) \mid (q, P_1, L_1, q') \in \Delta_1 \land L_1 \cap \SETX_2 = \emptyset \land r \in \{\bot, \top\}\} \\
		& \uplus \{((r, p), P_2, L_2, (r, p')) \mid (p, P_2, L_2, p') \in \Delta_2 \land L_2 \cap \SETX_1 = \emptyset \land r \in \{\bot, \top\}\} \\
		& \uplus \{((q, p), P_1 \land P_2, L_1 \cup L_2, (q', p')) \mid (q, P_1, L_1, q') \in \Delta_1 \land (p, P_2, L_2, p') \in \Delta_2 \land L_1 \cap \SETX_2 = L_2 \cap \SETX_1\} \\
		& \uplus \{((\top, \bot), \text{true}, \emptyset, (\top, \bot)), ((\bot, \top), \text{true}, \emptyset, (\bot, \top))\}
	\end{aligned}
	\]
	Intuitively, given a stream $\Stream$, $\cAjoins$ simulates $\cA_{\varphi_1}$ and $\cA_{\varphi_2}$ either independently or in parallel, ensuring at every step that any variable in the shared domain $\SETX_1 \cap \SETX_2$ is marked by both automata simultaneously ($L_1 \cap \SETX_2 = L_2 \cap \SETX_1$) and never marked when only one automaton is active.
	
	Next, we prove that the construction is correct, namely, $\sem{\varphi_1 \joins \varphi_2}(\Stream) = \sem{\cAjoins}(\Stream)$ for every stream $\Stream$. The proof is by double containment. 
	
	\paragraph{$\sem{\varphi_1 \joins \varphi_2}(\Stream) \subseteq \sem{\cAjoins}(\Stream)$}
	\vspace{4cm}
	
	\paragraph{$\sem{\cAjoins}(\Stream) \subseteq \sem{\varphi_1 \joins \varphi_2}(\Stream)$}
	\vspace{4cm}
	
\end{proof}